\section{Renormalized and jellium energies}
\label{sec:renormalized}

The universal-energy theorem is now proved. For the nonsummable Riesz
kernels with $0<s<2$ and for the logarithm, the meaningful finite comparison
comes from neutralizing the points against a uniform background and removing
their self-energy. Petrache and Serfaty established the implication from
universal optimality to these crystallization conclusions
\citep[Theorem~2]{PetracheSerfaty2020}. We give the transfer for the field
normalization and order of limits specified below.

The decisive intermediate comparison allows arbitrary period lattices.
Consequently the square-periodic approximation of the field infimum in
Appendix~\ref{app:normalized-periodic-approximation} suffices, although the
minimizing lattice is triangular.

\subsection{Field and scalar energy conventions}

Following the field-first construction of Petrache and Serfaty
\citep[Definitions~1.2--1.3]{PetracheSerfaty2017}, we define the energy
with its full-space normalization fixed by the fundamental-solution
equation below.  The truncation limit and cubic approximation for this
convention are established in
Proposition~\ref{prop:normalized-cubic-approximation}.
For $0\le s<2$, put
\[
 g_s(X)=\begin{cases}-\log|X|,&s=0,\\ |X|^{-s},&0<s<2,\end{cases}
 \qquad
 (k,\gamma)=\begin{cases}(0,0),&s=0,\\(1,s-1),&0<s<2.\end{cases}
\]
Write $X=(x,y)\in\R^2\times\R^k$, with the extra coordinate omitted
when $k=0$, and let $w_s=|y|^\gamma$ for $k=1$ and $w_0=1$.
The full-space identity
\[
 -\operatorname{div}(w_s\nabla g_s)=\kappa_s\delta_0
\]
gives $\kappa_0=2\pi$ and, for $0<s<2$,
\[
 \kappa_s=s\int_{S^2}|\omega_y|^{s-1}\,d\sigma(\omega)=4\pi.
\]
Here $\omega_y$ is the vertical coordinate on the unit sphere, and
$\delta_{\R^2}$ denotes area measure on $\R^2\times\{0\}$.

Fix once and for all $p$ with $1<p<2$ if $s=0$, and
\[
 1<p<\min\left(2,\frac2s,\frac3{s+1}\right)
 \qquad\text{if }0<s<2.
\]
For a locally finite configuration $\mathcal C$, admissible fields are
gradients $E=\nabla h$ in $L^p_{\mathrm{loc}}(\R^{2+k};\R^{2+k})$
with $w_sE\in L^1_{\mathrm{loc}}$, satisfying
\[
 -\operatorname{div}(w_sE)
 =\kappa_s\left(\sum_{p\in\mathcal C}\delta_{(p,0)}
                      -\delta_{\R^2}\right).
\]
This is the local integrability range following
Petrache and Serfaty~\citep[Definition~1.2]{PetracheSerfaty2017}.  Set
\[
 f_{s,\eta}=(g_s-g_s(\eta))_+,
 \qquad E_\eta=E-\sum_{p\in\mathcal C}
                       \nabla f_{s,\eta}(X-(p,0)).
\]
Thus the scalar potential subtracts $f_{s,\eta}$ and its gradient
subtracts $\nabla f_{s,\eta}$.  For $K_R=[-R/2,R/2]^2$, define
\begin{align*}
 \mathcal W_{s,\eta}(E)&=\limsup_{R\to\infty}
 \left[\frac1{R^2}\int_{K_R\times\R^k}w_s|E_\eta|^2
                         -\kappa_sg_s(\eta)\right],\\
 \mathcal W_s(E)&=\lim_{\eta\downarrow0}\mathcal W_{s,\eta}(E),
 \qquad W_s(\mathcal C)=\inf_E\mathcal W_s(E).
\end{align*}
The truncation limit exists in $\R\cup\{+\infty\}$ by
Proposition~\ref{prop:normalized-cubic-approximation}; the infimum is over
compatible admissible fields and is $+\infty$ if there are none.
In particular, the configuration infimum is outside both limits.

The periodic energy has the same field normalization.  For a torus
$T=\R^2/L$ of area $N\in\N$ and distinct points $a_1,\ldots,a_N$,
let $H$ be the canonical potential with mean-zero trace satisfying
\[
 -\operatorname{div}(w_s\nabla H)
 =\kappa_s\left(\sum_{i=1}^N\delta_{(a_i,0)}-\delta_T\right),
\]
with finite truncated energy on $T\times\R^k$.  Truncations include
the periodic images of the charges.  Set
\[
 W_{s,L}(a_1,\ldots,a_N)=\lim_{\eta\downarrow0}
 \left[\frac1N\int_{T\times\R^k}w_s|\nabla H_\eta|^2
                                      -\kappa_sg_s(\eta)\right].
\]
If $\mathcal G_L$ is the mean-zero unit-source extension Green function,
\[
 -\operatorname{div}(w_s\nabla\mathcal G_L)
      =\delta_{(0,0)}-N^{-1}\delta_T,
\]
integration by parts gives
\[
 W_{s,L}=\kappa_s^2\left[
  \frac1N\sum_{i\ne j}\mathcal G_L(a_i-a_j,0)
  +\lim_{x\to0}\left(\mathcal G_L(x,0)-\frac{g_s(x)}{\kappa_s}\right)
 \right].
\]
This follows the integration-by-parts argument of
Petrache and Serfaty~\citep[Proposition~1.5]{PetracheSerfaty2017}, using the Green function
and coefficient determined by the displayed extension equation.

For $N\in\N$, set $R=\sqrt N$ and $K_R=[-R/2,R/2]^2$.
The finite ordered-pair jellium energy is
\begin{equation}
\label{eq:finite-ordered-jellium}
\begin{split}
 J_{s,R}(a_1,\ldots,a_N)
 ={}&\sum_{i\ne j}g_s(a_i-a_j)
       -2\sum_{i=1}^N\int_{K_R}g_s(a_i-x)\,\dd x\\
    &+\int_{K_R}\int_{K_R}g_s(x-x')\,\dd x\,\dd x'.
\end{split}
\end{equation}
Here the particles are distinct; collisions have energy $+\infty$.
We use the density-one thermodynamic minimum
\[
 e_{\mathrm{Jel}}^{(s)}
 =\lim_{N\to\infty}\frac1N
       \min_{a_1,\ldots,a_N\in K_{\sqrt N}}
                    J_{s,\sqrt N}(a_1,\ldots,a_N).
\]
Since $|K_{\sqrt N}|=N$, normalization per particle and per unit area
coincide.  The scalar thermodynamic results and finite confinement
argument below establish the existence of this limit.

\subsection{The triangular minimum}

\begin{corollary}[Planar renormalized and jellium energies]
\label{cor:riesz-jellium}
Let $0\le s<2$.  For every positive integer $n$, put $N=n^2$ and
enumerate $A/(nA)$ by $a_1^\triangle,\ldots,a_N^\triangle$.
Every configuration of $N$ distinct points in $\R^2/(nA)$ satisfies
\[
 W_{s,nA}(a_1,\ldots,a_N)
 \ge W_{s,nA}(a_1^\triangle,\ldots,a_N^\triangle).
\]
Moreover, $A$ minimizes $W_s$ among density-one planar configurations,
and $e_{\mathrm{Jel}}^{(s)}=\kappa_s^{-1}W_s(A)$.
\end{corollary}

For a density-one periodic configuration $\mathcal C=\bigcup_{i=1}^N(a_i+L)$,
write $V_s(\mathcal C)=W_{s,L}(a_1,\ldots,a_N)$ for its canonical
periodic value. The following identity compares these values without
requiring the configurations to share a period lattice.

\begin{lemma}[Heat comparison across period lattices]
\label{lem:periodic-heat-comparison}
Let $0\le s<2$, and let $\mathcal C$ and $\mathcal D$ be simple periodic
configurations of density one. Put
\[
 q_0=2\pi,\qquad
 q_s=\frac{2^{2-s}\pi\Gamma(1-s/2)}{\Gamma(s/2)}\quad(0<s<2),
 \qquad \Psi_t(x)=(4\pi t)^{-1}e^{-|x|^2/(4t)}.
\]
If $\mathcal C=\bigcup_{i=1}^N(a_i+L)$, where $|\R^2/L|=N$, define
\[
 P_{t,L}(x)=\sum_{v\in L}\Psi_t(x-v),\qquad
 H_{\mathcal C}(t)=\frac1N\sum_{i,j=1}^N P_{t,L}(a_i-a_j)-\Psi_t(0),
\]
and define $H_{\mathcal D}$ using its own period lattice and point count.
Then the following integral is absolutely convergent and satisfies
\[
 V_s(\mathcal C)-V_s(\mathcal D)
 =\frac{\kappa_s q_s}{\Gamma(1-s/2)}\int_0^\infty
       \bigl(H_{\mathcal C}(t)-H_{\mathcal D}(t)\bigr)t^{-s/2}\,\dd t.
\]
\end{lemma}

\begin{proof}
Put $\alpha=1-s/2$. We first identify the scalar Green function whose
heat representation has the field normalization used here.
On the base plane, $(-\Delta)^\alpha g_s(\cdot,0)=q_s\delta_0$,
with the extra coordinate omitted when $s=0$.  On $T=\R^2/L$,
the spectral operator $(-\Delta)^\alpha$ multiplies the Fourier mode
$e^{2\pi i\xi\cdot x}$, $\xi\in L^*$, by $\lambda_\xi^\alpha$,
where $\lambda_\xi=(2\pi|\xi|)^2$.
Its mean-zero unit-source Green function $G_L$ therefore has coefficient
$N^{-1}\lambda_\xi^{-\alpha}$ for each $\xi\ne0$.
To compare it with the full-space extension, suppose first that $0<s<2$.
For $\lambda>0$, the even profile
\[
 u_\lambda(y)=\frac1{\Gamma(\alpha)}\int_0^\infty
 t^{\alpha-1}e^{-\lambda t-y^2/(4t)}\,\dd t
\]
decays at infinity and satisfies
\[
 u_\lambda(0)=\lambda^{-\alpha},\qquad
 u_\lambda''+\frac{1-2\alpha}{y}u_\lambda'-\lambda u_\lambda=0
 \quad(y>0).
\]
Differentiation and the substitution $v=y^2/(4t)$ give its full,
two-sided weighted flux:
\[
 -2\lim_{y\downarrow0}y^{1-2\alpha}u_\lambda'(y)
 =2^{2-2\alpha}\frac{\Gamma(1-\alpha)}{\Gamma(\alpha)}
 =\frac{\kappa_s}{q_s}.
\]
Thus the unit-source extension mode is
$(q_s/(N\kappa_s))u_{\lambda_\xi}$, whose trace is
$(q_s/\kappa_s)N^{-1}\lambda_\xi^{-\alpha}$.
The zero mode has no source: finite tail energy forces it to be
constant, and the mean-zero trace sets that constant to zero.
For $s=0$ there is no extension coordinate; both Green functions solve
$-\Delta G=\delta_0-N^{-1}$ on $T$ with zero mean, and
$q_0/\kappa_0=1$.  Hence, throughout $0\le s<2$,
$\mathcal G_L(x,0)=(q_s/\kappa_s)G_L(x)$, with the trace notation
omitted at $s=0$.  The periodic formula above consequently becomes
\[
 W_{s,L}=\kappa_sq_s\left[
  \frac1N\sum_{i\ne j}G_L(a_i-a_j)
  +\lim_{x\to0}\left(G_L(x)-\frac{g_s(x)}{q_s}\right)\right].
\]

Set
\[
 B_{\mathcal C}(t,x)=\frac1N\sum_{i,j}P_{t,L}(a_i-a_j+x),
\]
so $B_{\mathcal C}(t,0)=H_{\mathcal C}(t)+\Psi_t(0)$.
The heat representation \citep[Lemma~1]{PetracheSerfaty2020} is
\[
 G_L(x)=\frac1{\Gamma(\alpha)}\int_0^\infty
       \left(P_{t,L}(x)-\frac1N\right)t^{\alpha-1}\,\dd t.
\]
Its spectral proof applies to any flat torus.  To compare two periodic
configurations $\mathcal C$ and $\mathcal D$, set
$S_{\mathcal C}(x)=N^{-1}\sum_{i,j}G_L(a_i-a_j+x)$, and define
$S_{\mathcal D}$ using its own period lattice and point count.
Each $S$ has singular part $g_s(x)/q_s$.  In their heat-integrand
difference the background terms are both $1$ and the diagonal terms
are both $\Psi_t(x)$, so these terms cancel.  Uniformly for $x$ near
zero, the remaining difference is bounded by $Ct^{-1}e^{-c/t}$ as
$t\downarrow0$, because the other image distances are positive, and
by $Ce^{-\lambda t}$ as $t\to\infty$, by the two torus spectral gaps.
Dominated convergence as $x\to0$ consequently gives
\[
 V_s(\mathcal C)-V_s(\mathcal D)
 =\frac{\kappa_sq_s}{\Gamma(1-s/2)}\int_0^\infty
       \big(H_{\mathcal C}(t)-H_{\mathcal D}(t)\big)t^{-s/2}\,\dd t.
\]
This calculation also covers $s=0$: the logarithmic singularity
cancels before either endpoint is integrated.
\end{proof}

\begin{proof}[Proof of Corollary~\ref{cor:riesz-jellium}]
\emph{Periodic configurations.}
The subtraction in $H_{\mathcal C}$ removes only the coincident point;
all other images, including images of the same point, remain.
Periodic counting gives
\[
 H_{\mathcal C}(t)=\frac1{4\pi t}E_{g_t}(\mathcal C),
 \qquad g_t(u)=e^{-u/(4t)}.
\]
Indeed, every fixed image displacement has its periodic frequency in large
disks, and the Gaussian tail is summable uniformly over the finitely many
point types.  The Gaussian estimate \eqref{eq:gaussian-energy}, with
Gaussian parameter $1/(4\pi t)$, multiplied by the positive factor
$(4\pi t)^{-1}$, therefore gives $H_{\mathcal C}(t)\ge H_A(t)$
for every $t>0$.

Applying Lemma~\ref{lem:periodic-heat-comparison} with $\mathcal D=A$
therefore proves $V_s(\mathcal C)\ge V_s(A)$ for every periodic
$\mathcal C$. Repetition on $nA$ preserves $H_A$ and hence $V_s(A)$,
which proves the stated torus inequality.

\emph{The infinite-system field infimum.}

Put $m_s=\inf_E\mathcal W_s(E)$, where the infimum runs over all
compatible gradients in the fixed-$p$ class with simple locally finite
configurations and unit background. No centered-density assumption is
added to this infimum.
Proposition~\ref{prop:normalized-cubic-approximation} supplies canonical
periodic gradients whose values tend to $m_s$.  Each value is at least
$V_s(A)$ by the preceding inequality, so $m_s\ge V_s(A)$.
Conversely the canonical field of $A$
is admissible, so $m_s\le V_s(A)$.  For the configuration infimum we now
have
\[
 m_s\le W_s(A)\le V_s(A)=m_s,
\]
and every configuration has $W_s(\mathcal C)\ge m_s$. This lower bound
uses inclusion in the field class, not a conversion between disk and square
density. The cubic sequence suffices because the heat comparison applies to every
period lattice.

\emph{The scalar jellium value.}
Let $G_L$ again denote the mean-zero unit-source Green function of
$(-\Delta)^{1-s/2}$ on $\R^2/L$, as in the heat-comparison lemma.
To identify the jellium value, use the physical periodic kernel
$v_L=q_sG_L$ and its regular part
$r_L=\lim_{x\to0}(v_L(x)-g_s(x))$.  The unordered scalar periodic
energy, including its self term, is
\[
 \mathcal E_{\mathrm{per},L}
 =\sum_{i<j}v_L(a_i-a_j)+\frac N2r_L,
 \qquad
 \frac{2\mathcal E_{\mathrm{per},L}}N
 =\frac{W_{s,L}}{\kappa_s}.
\]
The scalar thermodynamic input, in this notation, is
\begin{equation}
\label{eq:scalar-jellium-input}
 \lim_{N\to\infty}\frac1{2N}
    \min_{\mathbf a\in(\R^2)^N}J_{s,\sqrt N}(\mathbf a)
 =
 \liminf_{N\to\infty}\frac1N
    \min_{\mathbf a\in(\R^2/(\sqrt N\Z^2))^N}
                   \mathcal E_{\mathrm{per},\sqrt N\Z^2}(\mathbf a)
 =:e_{\mathrm{per}}^{(s)} .
\end{equation}
For $0<s<2$, this is
Lewin et~al.~\citep[Section~II, (3)--(4), and Section~VI, (25), (27)--(31), Theorem~2]{LewinLiebSeiringer2019};
for $s=0$, it is
Lauritsen~\citep[Section~II and Theorem~II.1]{Lauritsen2021}, who also proves
existence of the full periodic limit.
Both inputs use $N$ particles against a background of density one
and area $N$, with unrestricted particle positions on the left.
Their scalar energies count each pair once.  Their periodic kernels
have zero mean, and their self contribution is $Nr_L/2$, exactly as
in $\mathcal E_{\mathrm{per},L}$ above.

The approximating periods in
Proposition~\ref{prop:normalized-cubic-approximation} are cubic.
Every canonical periodic value is at least $m_s$, and the values
along that sequence tend to $m_s$.  The finite identity
$2\mathcal E_{\mathrm{per},L}/N=W_{s,L}/\kappa_s$ therefore identifies
the lower limit in \eqref{eq:scalar-jellium-input} as
$e_{\mathrm{per}}^{(s)}=m_s/(2\kappa_s)$.

\emph{Confinement to the background square.}
Finite minimization is unchanged by restricting the particles to
$K_R$.  Fix the other $N-1$ particles.  For $0<s<2$, their one-particle
potential extends to $X=(x,z)\in\R^3$ as
\[
 U(X)=\sum_{j\ne i}|X-(a_j,0)|^{-s}
             -\int_{K_R}|X-(y,0)|^{-s}\,\dd y .
\]
The background potential is continuous because $s<2$, the point
poles tend to $+\infty$, and
$U(X)=-|X|^{-s}+O(|X|^{-s-1})$ at infinity.  Hence $U$ attains
a negative global minimum away from its point poles.

We use the weighted strong minimum principle:
a continuous weak solution
$u\in H^1_{\mathrm{loc}}(D,|z|^\gamma\,\dd X)$ of
$\operatorname{div}(|z|^\gamma\nabla u)=0$ on a connected open set
$D$, with $-1<\gamma<1$, cannot attain an interior minimum unless
it is constant
\citep[Corollary~2.3.10, p.~98]{FabesKenigSerapioni1982}.
This follows from their Harnack inequality for $A_2$ weights,
Theorem~2.3.8; here $|z|^\gamma\in A_2$.

Apply this principle on
\[
 D=\R^3\setminus\bigl((K_R\times\{0\})
                    \cup\{(a_j,0):j\ne i\}\bigr),
 \qquad \gamma=s-1 .
\]
The function $U$ is smooth away from these sources and belongs to
$H^1_{\mathrm{loc}}(D,|z|^{s-1}\,\dd X)$.
It is weighted harmonic also across the source-free part of $z=0$:
there $\partial_zU=O(|z|)$, so
$|z|^{s-1}\partial_zU=O(|z|^s)\to0$ from both sides.
The set $D$ is connected.  Constancy there is incompatible with
the negative minimum and $U\to0$ at infinity, so a minimum lies
on $K_R\times\{0\}$.
For $s=0$, ordinary harmonicity outside the square and the point
poles, together with $U(x)=\log|x|+o(1)$ at infinity, gives the same
conclusion by the ordinary strong minimum principle.
Moving each exterior particle to such a minimum does not increase
the energy or create a collision.  On $K_R^N$, the extended energy is
lower semicontinuous, finite at a distinct tuple, and diverges at
collisions, so its infimum is attained.  The confined and unrestricted
minima therefore agree, and
\[
 e_{\mathrm{Jel}}^{(s)}
 =2e_{\mathrm{per}}^{(s)}
 =\frac{m_s}{\kappa_s}.
\]
\end{proof}
